\repchapter{}{Головний висновок}
\begin{keybox}{Основна теза}
2 листопада 2026 року в Петропавлівській Борщагівці відбувся реальний локальний соціальний конфлікт, який переріс в інформаційну кризу загальноукраїнського масштабу. Його не можна звести ні до цілком вигаданої історії, ні до доведеної російської інформаційної операції.
\end{keybox}

Зібрані матеріали краще за інші моделі пояснює взаємодія трьох одночасних процесів: реального конфлікту навколо військовослужбовців та об'єкта Національної гвардії; інформаційної хвилі, у якій непідтверджені й емоційно сильні твердження випереджали перевірені відомості; і боротьби інтерпретацій, до якої долучилися українські паблики, державні установи, редакційні медіа, канали з попередніми російськими чи проросійськими мітками та звичайні користувачі.

Спостережувана картина не відповідає одній кампанії з доведеним центром управління. У досліджуваній панелі хвиля почалася з публікацій українських джерел, а «Украина.ру» долучилася о 13:57 з посиланням на київські ЗМІ. Найкраще підтверджена роль російського медіа в цьому епізоді --- подальше поширення сюжету. Доказів того, що російська держава створила інцидент, ініціювала першу хвилю чи керувала її учасниками, у зібраних матеріалах немає.

Окремий результат змінює спрощену двотаборову картину: 438 шаблонних коментарів від 367 акаунтів на підтримку українського посадовця мають статистично незвичну узгодженість. Можливу координацію потрібно оцінювати незалежно від політичної спрямованості повідомлень.

Лишаються невстановленими достовірність первинних тверджень про 40 незаконно утримуваних осіб, масові катування й загибель людини, час і причини травм чотирьох військовослужбовців, організатори конкретних дій та наявність зовнішнього керування. Розгорнуте обґрунтування подано в розділі~18.
